% Appendix: Discrete-sum check of the catalogue selection normalization
% Body only — \section{} and \label{sec:app:betag} are emitted by the builder.
% Primary sources: .planning/gate/G1_beta_g_check.md, .planning/gate/G1_beta_g_check.json
The \texttt{global} variant of the in-catalogue likelihood
(Sections~\ref{sec:pitfall} and~\ref{sec:estimators}) normalizes every
event by a single catalogue-wide selection sum rather than by a local,
per-event one. Its validity rests on an implicit bridge between a
discrete and a continuous object: the rate-weighted sum over the actual
GLADE+ catalogue \citep{Dalya:2021ewn},
\begin{equation}
  \Sigma_{\mathrm{glob}}(h) \;=\;
  \sum_{g\,:\,z_g < z_{\max}(h)} w_g\,
  p_{\mathrm{det}}\!\big(d_L(z_g, h)\big),
  \qquad
  w_g = \frac{R_{\mathrm{eff}}(M_g)}{1+z_g},
  \label{eq:app:sigmaglob}
\end{equation}
with $w_g$ the per-galaxy EMRI-rate weight of
Section~\ref{sec:framework} \citep{Babak:2017tow} and $z_{\max}(h)$ the
redshift reach of the population model, must track the continuous
in-catalogue selection integral $\beta_G(h) = D(h) - \beta_{\bar G}(h)$
up to an $h$-independent constant. If the catalogue were a
Monte-Carlo draw from a galaxy field of constant comoving number
density $\bar n_{\mathrm{gal}}$, one would have
\begin{equation}
  \Sigma_{\mathrm{glob}}(h) \;\approx\; C\,\beta_G(h),
  \qquad
  C = \bar n_{\mathrm{gal}}\,\big\langle R_{\mathrm{eff}} \big\rangle ,
  \label{eq:app:bridge}
\end{equation}
and the constant $C$ would cancel in the normalized posterior (see
Appendix~\ref{sec:app:gray} for the correspondence with the
catalogue-sum equations of \citealt{Gray:2019ksv}). Because this
cancellation is delicate, we test it by evaluating both sides of
Eq.~\eqref{eq:app:bridge} directly.

\medskip
\noindent\textit{Protocol.} Both sides are computed with identical
pipeline components: the injection-calibrated $p_{\mathrm{det}}$ built
from the seed600 injection pool with six sky bands at the nominal
threshold $\rho_{\mathrm{thr}} = 20$
% src: .planning/gate/G1_beta_g_check.json (n_sky_bands, snr_threshold)
(the pre-recalibration SNR scale; see Section~\ref{sec:realdata}), and
the per-pixel completeness $f_k$ entering $\beta_{\bar G}(h)$. The
discrete sum runs over the full reduced GLADE+ catalogue; the
continuous integrals $D(h)$ and $\beta_{\bar G}(h)$ are those of the
production likelihood. We evaluate the ratio
$\Sigma_{\mathrm{glob}}(h)/\beta_G(h)$ on $14$ values of $h$ in
$[0.60, 0.86]$ and normalize its shape at $h_{\mathrm{ref}} = 0.72$,
% src: .planning/gate/G1_beta_g_check.json (h grid, h_ref)
the grid point nearest the fiducial value. Only the $h$-shape of this
ratio is meaningful: any $h$-independent factor cancels, whereas an
$h$-dependent one multiplicatively distorts the in-catalogue channel.

\medskip
\noindent\textit{The expected $h^3$ factor.} The raw ratio is not
flat: it grows by a factor $2.48$ across the grid, an end-to-end tilt
of $+93$ per cent.
% src: .planning/gate/G1_beta_g_check.json (end_to_end_tilt_raw = 0.9285; ratio 8.957 -> 22.197)
The dominant contribution is expected. The catalogue is a
\emph{fixed} set of galaxies, so the comoving number density it
implies is not constant in $h$: distances scale as $1/h$, hence
$\bar n_{\mathrm{gal}}(h) \propto 1/V_c \propto h^3$, and
$(0.86/0.60)^3 = 2.94$ accounts for the bulk of the growth.
% src: .planning/gate/G1_beta_g_check.md (finding 1)
This factor is common to the numerator and denominator sums of
$L_{\mathrm{cat}}$ and cancels exactly in their ratio; it is therefore
not by itself a defect. It does, however, falsify the stronger
assumption that $\bar n_{\mathrm{gal}}$ is an overall constant: only
the \emph{ratio} cancellation is exact.

\medskip
\noindent\textit{The residual tilt.} After dividing out
$(h/h_{\mathrm{ref}})^3$, a smooth, monotonic residual remains
(Fig.~\ref{fig:betag}): the corrected shape falls from $1.085$ at
$h = 0.60$ to $0.913$ at $h = 0.86$, an end-to-end tilt of $-17.2$ per
cent ($+8.7$ to $-8.7$ per cent about $h_{\mathrm{ref}}$).
% src: .planning/gate/G1_beta_g_check.json (shape_h3_corrected endpoints; end_to_end_tilt_h3_corrected = -0.1719)
This is a genuine $h$-dependent mismatch between the discrete,
completeness-weighted content of the catalogue and the continuous
completeness-model integral, and unlike the $h^3$ factor it does
\emph{not} cancel: in \texttt{global} mode it multiplies the
in-catalogue channel $\beta_G L_{\mathrm{cat}}$ once per event, so the
tilt accumulates coherently over the $N \simeq 500$ events of the
sample
% src: .planning/gate/G1_beta_g_check.md (finding 3)
--- corroborating the railing behaviour of this mode seen in
Section~\ref{sec:realdata}. The local estimators are structurally
immune: \texttt{local\_ratio} and \texttt{volume\_deconv} never
evaluate $\Sigma_{\mathrm{glob}}$ (their $L_{\mathrm{cat}}$ is a local
self-normalized ratio in which ball-common factors cancel), and their
mixture weight $w_G = \beta_G/D$ is a ratio of two continuous
integrals.

\medskip
\noindent\textit{Origin and consequence.} Candidate origins of the
residual are (a)~the modelled completeness $f(z,\Omega,h)$ versus the
catalogue's true incompleteness as $z_{\max}(h)$ sweeps outward,
(b)~genuine inhomogeneity of the galaxy density (large-scale structure
and depth-dependent selection) versus the constant-comoving-density
assumption, and (c)~the surrogate of a mass-integrated rate constant
across the catalogue.
% src: .planning/gate/G1_beta_g_check.md (finding 2)
Disentangling these (e.g.\ by re-running with $f \equiv 1$ or on a
redshift-restricted subsample) is left to future work; for the present
analysis the decomposition in Fig.~\ref{fig:betag} suffices. The raw
factor $\sim\!2.5$ that cancels and the residual $-17$ per cent that
does not are, together, the anatomy of why global-denominator
normalizations of galaxy-catalogue dark-siren likelihoods are fragile:
we accordingly exclude the \texttt{global} mode from all quoted
results.
